\documentclass[letterpaper]{article}
\usepackage[submission]{aaai2027}
\usepackage[hyphens]{url}
\usepackage{graphicx}
\urlstyle{rm}
\def\UrlFont{\rm}
\usepackage{natbib}
\usepackage{caption}
\usepackage{booktabs}
\usepackage{amsmath}
\frenchspacing
\pdfinfo{
/TemplateVersion (2027.1)
/Title (LeakBench-Tab Technical Supplement)
/Author (Anonymous Submission)
/Subject (AAAI-27 Supplementary Material)
}

% Safe defaults used until generated/result_macros.tex is produced from the
% complete confirmatory paper_claims.json.  The front-page draft banner is
% controlled by \ifLBResultsReady in main.tex and supplement.tex.
\newif\ifLBResultsReady
\newif\ifLBSimpleStructuredSupported
\LBResultsReadyfalse
\LBSimpleStructuredSupportedfalse

\newcommand{\LBPending}{\ensuremath{\mathrm{PENDING}}}
\newcommand{\LBSourceSHA}{\LBPending}
\newcommand{\LBCDXScatterPath}{figures/generated/cdx_scatter.pdf}
\newcommand{\LBMechanismModelHeatmapPath}{figures/generated/mechanism_model_heatmap.pdf}
\newcommand{\LBStrengthDiagnosticPath}{figures/generated/strength_diagnostic_robustness.pdf}
\newcommand{\LBControlledTaskCount}{20}
\newcommand{\LBMechanismCount}{11}
\newcommand{\LBStrengthCount}{5}
\newcommand{\LBCoreModelCount}{5}
\newcommand{\LBSeedCount}{5}
\newcommand{\LBExpectedCells}{27,500}
\newcommand{\LBSuccessfulCells}{\LBPending}
\newcommand{\LBCompletionRate}{\LBPending}
\newcommand{\LBDiagnosticMethodCount}{4}
\newcommand{\LBExpectedDiagnosticCells}{22,000}
\newcommand{\LBSuccessfulDiagnosticCells}{\LBPending}
\newcommand{\LBDiagnosticCompletionRate}{\LBPending}
\newcommand{\LBDiagnosticConditionalStatus}{pending}
\newcommand{\LBSimpleStructuredStatus}{pending}
\newcommand{\LBSimpleStructuredDifference}{\LBPending}
\newcommand{\LBSimpleStructuredCILow}{\LBPending}
\newcommand{\LBSimpleStructuredCIHigh}{\LBPending}
\newcommand{\LBSimpleStructuredHolmP}{\LBPending}
\newcommand{\LBModelPositiveDirectionCount}{\LBPending}
\newcommand{\LBModelCIExcludesZeroCount}{\LBPending}

\newcommand{\LBMThreeStatus}{pending}
\newcommand{\LBMThreeDetectability}{\LBPending}
\newcommand{\LBMThreeDetectabilityCILow}{\LBPending}
\newcommand{\LBMThreeDetectabilityCIHigh}{\LBPending}
\newcommand{\LBMThreeHarm}{\LBPending}
\newcommand{\LBMThreeHarmCILow}{\LBPending}
\newcommand{\LBMThreeHarmCIHigh}{\LBPending}

\newcommand{\LBMEightStatus}{pending}
\newcommand{\LBMEightDetectability}{\LBPending}
\newcommand{\LBMEightDetectabilityCILow}{\LBPending}
\newcommand{\LBMEightDetectabilityCIHigh}{\LBPending}
\newcommand{\LBMEightHarm}{\LBPending}
\newcommand{\LBMEightHarmCILow}{\LBPending}
\newcommand{\LBMEightHarmCIHigh}{\LBPending}
\newcommand{\LBMEightClusterCILow}{\LBPending}
\newcommand{\LBMEightClusterCIHigh}{\LBPending}

\newcommand{\LBMNineStatus}{pending}
\newcommand{\LBMNineDetectability}{\LBPending}
\newcommand{\LBMNineDetectabilityCILow}{\LBPending}
\newcommand{\LBMNineDetectabilityCIHigh}{\LBPending}
\newcommand{\LBMNineHarm}{\LBPending}
\newcommand{\LBMNineHarmCILow}{\LBPending}
\newcommand{\LBMNineHarmCIHigh}{\LBPending}
\newcommand{\LBMNineReweightingLow}{\LBPending}
\newcommand{\LBMNineReweightingHigh}{\LBPending}

\newcommand{\LBRelationStatus}{pending}
\newcommand{\LBGlobalSpearman}{\LBPending}
\newcommand{\LBGlobalSpearmanCILow}{\LBPending}
\newcommand{\LBGlobalSpearmanCIHigh}{\LBPending}
\newcommand{\LBCategoryRSquared}{\LBPending}
\newcommand{\LBCategoryPlusDRSquared}{\LBPending}
\newcommand{\LBIncrementalRSquared}{\LBPending}
\newcommand{\LBIncrementalPermutationP}{\LBPending}
\newcommand{\LBCategoryLomoRSquared}{\LBPending}
\newcommand{\LBCategoryPlusDLomoRSquared}{\LBPending}
\newcommand{\LBIncrementalLomoRSquared}{\LBPending}

\newcommand{\LBNaturalStatus}{pending}
\newcommand{\LBNaturalTaskCount}{5}
\newcommand{\LBNaturalModelCount}{4}
\newcommand{\LBNaturalAllPositive}{\LBPending}
\newcommand{\LBNaturalMeanHarm}{\LBPending}
\newcommand{\LBNaturalCILow}{\LBPending}
\newcommand{\LBNaturalCIHigh}{\LBPending}
\newcommand{\LBNaturalSignFlipP}{\LBPending}

\newcommand{\LBMThreeMIDetectability}{\LBPending}
\newcommand{\LBMThreeMICILow}{\LBPending}
\newcommand{\LBMThreeMICIHigh}{\LBPending}
\newcommand{\LBMThreeCorrelationDetectability}{\LBPending}
\newcommand{\LBMThreeCorrelationCILow}{\LBPending}
\newcommand{\LBMThreeCorrelationCIHigh}{\LBPending}
\newcommand{\LBMThreeLRCoefficientDetectability}{\LBPending}
\newcommand{\LBMThreeLRCoefficientCILow}{\LBPending}
\newcommand{\LBMThreeLRCoefficientCIHigh}{\LBPending}
\newcommand{\LBMThreeRFPermutationDetectability}{\LBPending}
\newcommand{\LBMThreeRFPermutationCILow}{\LBPending}
\newcommand{\LBMThreeRFPermutationCIHigh}{\LBPending}
\newcommand{\LBMThreeDiagnosticMin}{\LBPending}
\newcommand{\LBMThreeDiagnosticMax}{\LBPending}

\newcommand{\LBMFourMIDetectability}{\LBPending}
\newcommand{\LBMFourMICILow}{\LBPending}
\newcommand{\LBMFourMICIHigh}{\LBPending}
\newcommand{\LBMFourCorrelationDetectability}{\LBPending}
\newcommand{\LBMFourCorrelationCILow}{\LBPending}
\newcommand{\LBMFourCorrelationCIHigh}{\LBPending}
\newcommand{\LBMFourLRCoefficientDetectability}{\LBPending}
\newcommand{\LBMFourLRCoefficientCILow}{\LBPending}
\newcommand{\LBMFourLRCoefficientCIHigh}{\LBPending}
\newcommand{\LBMFourRFPermutationDetectability}{\LBPending}
\newcommand{\LBMFourRFPermutationCILow}{\LBPending}
\newcommand{\LBMFourRFPermutationCIHigh}{\LBPending}
\newcommand{\LBMFourDiagnosticMin}{\LBPending}
\newcommand{\LBMFourDiagnosticMax}{\LBPending}

\newcommand{\LBMFiveMIDetectability}{\LBPending}
\newcommand{\LBMFiveMICILow}{\LBPending}
\newcommand{\LBMFiveMICIHigh}{\LBPending}
\newcommand{\LBMFiveCorrelationDetectability}{\LBPending}
\newcommand{\LBMFiveCorrelationCILow}{\LBPending}
\newcommand{\LBMFiveCorrelationCIHigh}{\LBPending}
\newcommand{\LBMFiveLRCoefficientDetectability}{\LBPending}
\newcommand{\LBMFiveLRCoefficientCILow}{\LBPending}
\newcommand{\LBMFiveLRCoefficientCIHigh}{\LBPending}
\newcommand{\LBMFiveRFPermutationDetectability}{\LBPending}
\newcommand{\LBMFiveRFPermutationCILow}{\LBPending}
\newcommand{\LBMFiveRFPermutationCIHigh}{\LBPending}
\newcommand{\LBMFiveDiagnosticMin}{\LBPending}
\newcommand{\LBMFiveDiagnosticMax}{\LBPending}

\IfFileExists{generated/result_macros.tex}{% AUTO-GENERATED by source_data/generate_result_macros.py.
% Do not edit by hand.
% paper_claims.json sha256: deff6fe15791439d16c64673061091eed62b5fc4f0f6bc60780fe7b5dfb64c04
\LBResultsReadytrue
\LBSimpleStructuredSupportedtrue
\renewcommand{\LBSourceSHA}{deff6fe15791439d16c64673061091eed62b5fc4f0f6bc60780fe7b5dfb64c04}
\renewcommand{\LBControlledTaskCount}{20}
\renewcommand{\LBMechanismCount}{11}
\renewcommand{\LBStrengthCount}{5}
\renewcommand{\LBCoreModelCount}{5}
\renewcommand{\LBSeedCount}{5}
\renewcommand{\LBExpectedCells}{27,500}
\renewcommand{\LBSuccessfulCells}{27,500}
\renewcommand{\LBCompletionRate}{100.0}
\renewcommand{\LBDiagnosticMethodCount}{4}
\renewcommand{\LBExpectedDiagnosticCells}{22,000}
\renewcommand{\LBSuccessfulDiagnosticCells}{22,000}
\renewcommand{\LBDiagnosticCompletionRate}{100.0}
\renewcommand{\LBDiagnosticConditionalStatus}{descriptive only}
\renewcommand{\LBSimpleStructuredStatus}{supported}
\renewcommand{\LBSimpleStructuredDifference}{0.164}
\renewcommand{\LBSimpleStructuredCILow}{0.150}
\renewcommand{\LBSimpleStructuredCIHigh}{0.178}
\renewcommand{\LBSimpleStructuredHolmP}{<0.001}
\renewcommand{\LBModelPositiveDirectionCount}{5}
\renewcommand{\LBModelCIExcludesZeroCount}{0}
\renewcommand{\LBMThreeStatus}{descriptive only}
\renewcommand{\LBMThreeDetectability}{0.079}
\renewcommand{\LBMThreeDetectabilityCILow}{0.045}
\renewcommand{\LBMThreeDetectabilityCIHigh}{0.128}
\renewcommand{\LBMThreeHarm}{0.217}
\renewcommand{\LBMThreeHarmCILow}{0.201}
\renewcommand{\LBMThreeHarmCIHigh}{0.234}
\renewcommand{\LBMEightStatus}{descriptive only}
\renewcommand{\LBMEightDetectability}{0.356}
\renewcommand{\LBMEightDetectabilityCILow}{0.263}
\renewcommand{\LBMEightDetectabilityCIHigh}{0.457}
\renewcommand{\LBMEightHarm}{0.003}
\renewcommand{\LBMEightHarmCILow}{-0.002}
\renewcommand{\LBMEightHarmCIHigh}{0.008}
\renewcommand{\LBMNineStatus}{descriptive only}
\renewcommand{\LBMNineDetectability}{0.541}
\renewcommand{\LBMNineDetectabilityCILow}{0.500}
\renewcommand{\LBMNineDetectabilityCIHigh}{0.583}
\renewcommand{\LBMNineHarm}{0.235}
\renewcommand{\LBMNineHarmCILow}{0.216}
\renewcommand{\LBMNineHarmCIHigh}{0.254}
\renewcommand{\LBMEightClusterCILow}{-0.006}
\renewcommand{\LBMEightClusterCIHigh}{0.010}
\renewcommand{\LBMNineReweightingLow}{0.199}
\renewcommand{\LBMNineReweightingHigh}{0.235}
\renewcommand{\LBRelationStatus}{descriptive only}
\renewcommand{\LBGlobalSpearman}{0.555}
\renewcommand{\LBGlobalSpearmanCILow}{0.482}
\renewcommand{\LBGlobalSpearmanCIHigh}{0.709}
\renewcommand{\LBCategoryRSquared}{0.514}
\renewcommand{\LBCategoryPlusDRSquared}{0.586}
\renewcommand{\LBIncrementalRSquared}{0.072}
\renewcommand{\LBIncrementalPermutationP}{0.330}
\renewcommand{\LBCategoryLomoRSquared}{0.086}
\renewcommand{\LBCategoryPlusDLomoRSquared}{-0.993}
\renewcommand{\LBIncrementalLomoRSquared}{-1.079}
\renewcommand{\LBNaturalStatus}{case-study only}
\renewcommand{\LBNaturalTaskCount}{5}
\renewcommand{\LBNaturalModelCount}{4}
\renewcommand{\LBNaturalAllPositive}{yes}
\renewcommand{\LBNaturalMeanHarm}{0.273}
\renewcommand{\LBNaturalCILow}{0.120}
\renewcommand{\LBNaturalCIHigh}{0.424}
\renewcommand{\LBNaturalSignFlipP}{0.062}
\renewcommand{\LBMThreeMIDetectability}{0.079}
\renewcommand{\LBMThreeMICILow}{0.069}
\renewcommand{\LBMThreeMICIHigh}{0.089}
\renewcommand{\LBMThreeCorrelationDetectability}{0.035}
\renewcommand{\LBMThreeCorrelationCILow}{0.025}
\renewcommand{\LBMThreeCorrelationCIHigh}{0.045}
\renewcommand{\LBMThreeLRCoefficientDetectability}{0.045}
\renewcommand{\LBMThreeLRCoefficientCILow}{0.035}
\renewcommand{\LBMThreeLRCoefficientCIHigh}{0.055}
\renewcommand{\LBMThreeRFPermutationDetectability}{0.602}
\renewcommand{\LBMThreeRFPermutationCILow}{0.592}
\renewcommand{\LBMThreeRFPermutationCIHigh}{0.612}
\renewcommand{\LBMThreeDiagnosticMin}{0.035}
\renewcommand{\LBMThreeDiagnosticMax}{0.602}
\renewcommand{\LBMFourMIDetectability}{0.081}
\renewcommand{\LBMFourMICILow}{0.071}
\renewcommand{\LBMFourMICIHigh}{0.091}
\renewcommand{\LBMFourCorrelationDetectability}{0.039}
\renewcommand{\LBMFourCorrelationCILow}{0.029}
\renewcommand{\LBMFourCorrelationCIHigh}{0.049}
\renewcommand{\LBMFourLRCoefficientDetectability}{0.037}
\renewcommand{\LBMFourLRCoefficientCILow}{0.027}
\renewcommand{\LBMFourLRCoefficientCIHigh}{0.047}
\renewcommand{\LBMFourRFPermutationDetectability}{0.045}
\renewcommand{\LBMFourRFPermutationCILow}{0.035}
\renewcommand{\LBMFourRFPermutationCIHigh}{0.055}
\renewcommand{\LBMFourDiagnosticMin}{0.037}
\renewcommand{\LBMFourDiagnosticMax}{0.081}
\renewcommand{\LBMFiveMIDetectability}{0.084}
\renewcommand{\LBMFiveMICILow}{0.074}
\renewcommand{\LBMFiveMICIHigh}{0.094}
\renewcommand{\LBMFiveCorrelationDetectability}{0.037}
\renewcommand{\LBMFiveCorrelationCILow}{0.027}
\renewcommand{\LBMFiveCorrelationCIHigh}{0.047}
\renewcommand{\LBMFiveLRCoefficientDetectability}{0.036}
\renewcommand{\LBMFiveLRCoefficientCILow}{0.026}
\renewcommand{\LBMFiveLRCoefficientCIHigh}{0.046}
\renewcommand{\LBMFiveRFPermutationDetectability}{0.044}
\renewcommand{\LBMFiveRFPermutationCILow}{0.034}
\renewcommand{\LBMFiveRFPermutationCIHigh}{0.054}
\renewcommand{\LBMFiveDiagnosticMin}{0.036}
\renewcommand{\LBMFiveDiagnosticMax}{0.084}
}{}

\title{LeakBench-Tab: Technical Supplement}
\author{Anonymous Submission}
\affiliations{}

\begin{document}
\maketitle

\ifLBResultsReady\else
\begin{center}
\fbox{\parbox{0.94\linewidth}{\centering\bfseries
INTERNAL DRAFT --- RESULTS BLOCKED.\\
This supplement contains no submission-ready empirical result until the final
confirmatory claim release generates the evidence macros.}}
\end{center}
\fi

\section{Scope and Artifact Contract}
This supplement covers the corrected controlled benchmark and five fixed
real-data case studies.  The controlled design contains
\LBControlledTaskCount{} tasks, \LBMechanismCount{} mechanisms,
\LBStrengthCount{} strengths, \LBCoreModelCount{} model families, and
\LBSeedCount{} injection seeds (\LBExpectedCells{} planned cells).  Metadata
localization and governance experiments are not part of the confirmatory claim
set and are therefore not reported as paper findings.  A separate
\LBExpectedDiagnosticCells{}-cell matrix compares
\LBDiagnosticMethodCount{} oracle-blind, pre-test development rankers on the
same frozen tasks; diagnostic rows are not counted as model-training cells.

Numerical macros are generated only from
\texttt{results/corrected\_v2/paper\_claims.json}.  The generator verifies a
complete confirmatory matrix, final claim statuses, finite ordered intervals,
the five-model identity set, all \LBExpectedDiagnosticCells{} diagnostic cells,
and evidence provenance before writing LaTeX.  Its source digest is
\texttt{\LBSourceSHA}.

\section{Mechanism Contracts and Construction Invariants}
Every injected task records a contaminated-field mask, timestamps, entity and
source identifiers, feature-availability metadata, structural groups, graph
edges, and immutable task and split hashes.  C is fixed by these construction
contracts before diagnostic or model fitting.

\begin{table*}[t]
\centering
\small
\begin{tabular}{llll}
\toprule
IDs & Category & Required invariant & Construction audit \\
\midrule
M01, M02 & Simple & Target-derived field; strength changes signal/noise & Separation by strength \\
M06 & Simple & Redundant fields form one declared contaminated group & Group size matches level \\
M10 & Simple & Legitimate component equals clean input exactly & Equality plus separate mask \\
M03 & Boundary & Nonlinear target interaction & Interaction and label audit \\
M07 & Boundary & Contamination restricted to declared subgroup & Prevalence audit \\
M11 & Boundary & Contaminated fields form declared graph component & Component audit \\
M04, M05 & Structured & Current and past rows excluded from future statistic & Temporal index audit \\
M08 & Structured & Future entity outcome; current label excluded & Leave-one-out/prior audit \\
M09 & Structured & Outcome-dependent source collection; complete one-hot & Source-ID permutation invariance \\
\bottomrule
\end{tabular}
\caption{Construction invariants that must hold independently of empirical
effect size.  Final release tables should cite the corresponding automated
audit and frozen task-manifest hash.}
\label{tab:construction-invariants}
\end{table*}

\subsection{Strength Mapping}
The five generic strength labels map to $\{0.2,0.4,0.6,0.8,1.0\}$.  Mechanisms
with discrete structural strength use frozen mechanism-specific mappings:
M06 redundant-field counts $\{1,2,4,6,8\}$, M07 subgroup prevalences
$\{0.05,0.10,0.20,0.35,0.50\}$, and M11 component sizes
$\{2,3,4,6,8\}$.  M04, M05, M08, and M09 must be documented with their exact
future-window, shrinkage, support, and source-assignment parameters from the
frozen corrected protocol rather than reconstructed from result tables.
M09 represents one semantic source variable with eight complete one-hot
columns.  Its AP evaluates encoded-column localization against those eight
oracle-marked columns; M09 D is therefore representation-conditional and is
not evidence that a source field is equally detectable under arbitrary
encodings or field-level scoring rules.

M04 and M05 use row-count windows after stable timestamp sorting, not fixed
calendar durations.  For a row at time $t$, both begin strictly after every row
with timestamp $t$ (a right-sided search), so the current row and all
same-timestamp rows are excluded.  M04 then averages the next
$\max(2,\lfloor0.05n\rfloor)$ rows.  M05 first skips
$\max(1,\lfloor n\delta\rfloor)-1$ strictly future rows and averages the next
$\max(2,\max(1,\lfloor n\delta\rfloor))$ rows, with the frozen
$\delta=0.02$ at every strength.  A future timestamp tie can be truncated by
the row-count width.  Consequently these mechanisms are sensitive to sampling
density and stable order within a future tie; they represent controlled
row-window contamination rather than an invariant duration-based event-time
estimand.  If a row has no eligible strictly future observations after the
declared skip, the constructed future statistic falls back to the leave-one-out
global label prior; the current row's label is never reintroduced by that
fallback.

\section{Controlled Task Registry and Splits}
The task registry contains twenty deterministic panel generators across five
archetypes.  Generator randomness is independent of injection randomness.
Rows are ordered chronologically and assigned once to 60\% train, 20\%
validation, and 20\% test partitions.  The split is invariant to mechanism,
strength, and downstream learner.  The final camera-ready appendix should list
for each task: task identifier, archetype, row and feature counts, entity and
source support, class prevalence, generator seed, and split hash.  These values
must be imported from the canonical manifest, not manually transcribed.

\section{Model Identity and Preprocessing}
The core set comprises logistic regression, random forest
\citep{breiman2001random}, CatBoost \citep{prokhorenkova2018catboost}, LightGBM
\citep{ke2017lightgbm}, and the official TabM implementation
\citep{gorishniy2025tabm}.  All scaling and learned preprocessing
use training rows only.  Strict and permissive views share row partitions,
labels, optimization seeds, and model hyperparameters.  The official TabM run
uses immutable serialized task bundles and verifies bundle, task, and split
hashes before fitting.  The release must include package versions, CUDA and GPU
identity, determinism settings, optimizer, learning rate, batch size, epoch and
early-stopping rules, and per-model preprocessing.  Surrogate implementations
and silent fallbacks are prohibited.

\begin{table*}[t]
\centering
\scriptsize
\setlength{\tabcolsep}{3pt}
\begin{tabular}{@{}p{0.10\textwidth}p{0.17\textwidth}p{0.34\textwidth}p{0.31\textwidth}@{}}
\toprule
Learner & Train-fitted preprocessing & Frozen fit configuration & Selection / prediction \\
\midrule
LR & StandardScaler & $C=1$, max 2,000 iterations & Test positive-class probability \\
RF & None & 250 trees, min leaf 2, max features $\sqrt{p}$ & Test positive-class probability \\
CatBoost & None & 250 iterations, rate 0.05, depth 6, CPU & Validation AUC; patience 30 \\
LightGBM & None & 250 trees, rate 0.05, 31 leaves, depth 6, min child 20 &
Validation binary log loss; patience 30 \\
TabM & StandardScaler & Official TabM v0.0.3, $k=32$, AdamW, rate 0.003,
weight decay $10^{-4}$ & Validation BCE; max 200 epochs, patience 20, batch
256; mean of per-$k$ sigmoid probabilities \\
\bottomrule
\end{tabular}
\caption{Frozen learner configurations.  Every strict/permissive pair shares
the split, seed, validation rule, and learner configuration.  CPU learners use
one thread; official TabM uses deterministic PyTorch algorithms on the declared
CUDA device.}
\label{tab:model-configurations}
\end{table*}

\section{Estimands and Statistical Procedures}
For a downstream learner $m$, the paired effect is
\begin{equation}
X_m = \mathrm{AUROC}_{m,\mathrm{permissive}}
      -\mathrm{AUROC}_{m,\mathrm{strict}}.
\end{equation}
For a feature-localization diagnostic, normalized average precision is
\begin{equation}
D = \frac{\mathrm{AP}-\pi}{1-\pi},
\end{equation}
where $\pi$ is contaminated-column prevalence.  Mutual-information scores are
computed from training rows once per injected task with fixed
\texttt{random\_state=42} and must be identical across all five downstream
learners.  Three additional oracle-blind rankers---absolute
correlation, logistic coefficient magnitude, and random-forest permutation
importance---form a frozen sensitivity suite.  MI, absolute correlation, and
logistic coefficients use training rows only.  The random forest is fitted on
training rows, but permutation importance is evaluated against frozen
validation labels.  Thus all four are pre-test development diagnostics, not
four training-only diagnostics.  D is always labeled by diagnostic and encoding; a
per-mechanism maximum is not a deployable selector.

The independent unit is the controlled task.  The 20,000-draw hierarchical
bootstrap first samples tasks with replacement and then samples injection seeds
within each selected task.  Task-level sign flips test mechanism effects, with
Holm adjustment over the eleven-mechanism family.  The three declared category
contrasts form a separate Holm family.  Model-specific simple-minus-structured
contrasts form a five-model family.  The category tests are exact two-sided
sign flips over the 20 task effects, not bootstrap-tail $p$-values.  M08 samples
one entity vector per dataset and applies it jointly to all 125
seed-by-model-by-strength cells.  It retains a separate effect for each
injection seed under that shared entity draw before outer dataset/seed
resampling.  M09 samples one
source-category vector per dataset/seed/strength and applies it jointly to all
five models.  Its eight source IDs are designed categories rather than draws
from a source population, so the resulting interval is descriptive reweighting
only.  All failed cells remain in the raw ledger; incomplete matrices are not
analyzed as confirmatory.

The diagnostic sensitivity matrix crosses twenty tasks, eleven mechanisms,
five strengths, five seeds, and four rankers, for 22,000 cells.  The first
three rankers use training rows only; RF permutation trains on training rows
and evaluates perturbations on the frozen validation split.  Intervals use the
same task/seed hierarchy as the primary D analysis.  Method comparisons are
descriptive sensitivity analyses; no paired simultaneous method-comparison
interval was frozen, and no ranker is chosen after observing contamination
labels.

\paragraph{Frozen diagnostic RNG amendment.}
After the confirmatory sensitivity suite was complete, an evidence-chain audit
found a provenance mismatch: its mutual-information implementation used the
injection seed, whereas the frozen headline implementation and the
immutable task manifest used fixed seed 42.  The post-unblinding amendment is
deterministic and scope-locked: it replaces the AP, normalized-AP, and top-5
recall fields on every and only the 5,500 mutual-information rows with the
identity-matched values already stored in the frozen task manifest.  It
preserves every original field on the other 16,500 rows and does not change
the four-method set.  The amendment builder reads no
model-outcome columns and performs no value-dependent row selection or tuning.
The task-seeded raw suite is retained only for provenance; the paper and
release validator consume the amended canonical diagnostic table.  All
diagnostic-method and mechanism-profile comparisons are descriptive; the
amendment does not authorize thresholded profile conclusions.

\paragraph{Frozen statistical inference amendment.}
A separate post-unblinding methods audit identified three invalid uncertainty
constructions in the superseded outputs: category bootstrap tail areas had been
labeled as $p$-values, D--X intervals resampled X while holding D fixed, and
cluster draws were independent across cells generated from the same task.  A
second audit found that the initial cluster amendment still drew entities
independently across M08 injection seeds that reuse the same test rows, labels,
and entities.  Before final amended analysis, the second replacement code and
immutable task registry were hash-frozen.  The replacement uses exact task-level sign flips with Holm
adjustment for the three category contrasts, jointly resamples D and X with the
same dataset/seed indices, and synchronizes entity/source-category draws at the
grouping levels above.  M03, M08, M09, D--X, diagnostic-method, and model-specific
summaries remain descriptive; only the pre-existing simple-versus-structured
directional contrast uses the exact Holm-and-CI support rule.  Output manifests
bind analysis code, canonical data, configuration, all twenty frozen task
bundles, and every consumed M08/M09 prediction-file hash; row IDs, labels,
entities, sources, clean/full AUC, and task hashes are checked directly.  The
release validator deep-reruns every paper-facing statistical table.

After all protocol, amendment, support-rule, and claim-scope freezes, a local
hyperparameter search unintentionally displayed a truncated subset of
in-progress GPU checkpoint rows.  The incident is retained in the artifact
audit trail; it authorized no task exclusion, mechanism change, inferential
change, threshold change, or manuscript result wording.  All reported values
are derived only after completion from the fully validated canonical release.

\section{Complete Controlled Results}
\ifLBResultsReady
The canonical matrix contains \LBSuccessfulCells{} successful cells
(\LBCompletionRate\%).  The final supplement must include the complete
mechanism-by-model point estimates, hierarchical intervals, adjusted
$p$-values, and strength-response curves generated from the canonical table.

\begin{table}[t]
\centering
\small
\begin{tabular}{lrrr}
\toprule
Profile & Normalized AP & Paired harm & Claim status \\
\midrule
M03 & \LBMThreeDetectability{} & \LBMThreeHarm{} & \LBMThreeStatus{} \\
M08 & \LBMEightDetectability{} & \LBMEightHarm{} & \LBMEightStatus{} \\
M09 & \LBMNineDetectability{} & \LBMNineHarm{} & \LBMNineStatus{} \\
\bottomrule
\end{tabular}
\caption{Headline crossed profiles.  Full intervals and per-model estimates
belong in the machine-generated extended table.}
\label{tab:headline-profiles}
\end{table}

The simple-minus-structured contrast is
$\LBSimpleStructuredDifference$ (95\% CI
$[\LBSimpleStructuredCILow,\LBSimpleStructuredCIHigh]$; Holm-adjusted
$p=\LBSimpleStructuredHolmP$ from an exact task-level sign-flip test).  It is positive in
\LBModelPositiveDirectionCount{}/5 model families, with
\LBModelCIExcludesZeroCount{}/5 model-specific intervals excluding zero.
\else
Complete controlled results are intentionally absent from this draft.  No
pilot estimate is a substitute for the final canonical matrix.
\fi

\section{Diagnostic-Conditional Detectability}
\ifLBResultsReady
The diagnostic matrix contains \LBSuccessfulDiagnosticCells{}/
\LBExpectedDiagnosticCells{} successful cells
(\LBDiagnosticCompletionRate\%).  Table~\ref{tab:diagnostic-sensitivity} reports
the three mechanisms targeted by the frozen sensitivity claim; complete
eleven-mechanism tables and confidence intervals belong in the archival
supplement.

\begin{table*}[t]
\centering
\small
\begin{tabular}{lrrrrr}
\toprule
Mechanism & Mutual information & Absolute correlation & LR coefficient & RF permutation & Range \\
\midrule
M03 & \LBMThreeMIDetectability{} & \LBMThreeCorrelationDetectability{} &
\LBMThreeLRCoefficientDetectability{} & \LBMThreeRFPermutationDetectability{} &
[\LBMThreeDiagnosticMin{}, \LBMThreeDiagnosticMax{}] \\
M04 & \LBMFourMIDetectability{} & \LBMFourCorrelationDetectability{} &
\LBMFourLRCoefficientDetectability{} & \LBMFourRFPermutationDetectability{} &
[\LBMFourDiagnosticMin{}, \LBMFourDiagnosticMax{}] \\
M05 & \LBMFiveMIDetectability{} & \LBMFiveCorrelationDetectability{} &
\LBMFiveLRCoefficientDetectability{} & \LBMFiveRFPermutationDetectability{} &
[\LBMFiveDiagnosticMin{}, \LBMFiveDiagnosticMax{}] \\
\bottomrule
\end{tabular}
\caption{Normalized AP under four frozen oracle-blind, pre-test development
rankers.  The first three use training rows only; RF permutation uses frozen
validation labels.  Mutual information is primary.  The comparison status is
\emph{\LBDiagnosticConditionalStatus{}}; the row-wise maximum is not an
operational method-selection result.}
\label{tab:diagnostic-sensitivity}
\end{table*}
\else
Diagnostic sensitivity values remain blocked.  In particular, no pilot
method-ordering is transcribed into the manuscript.
\fi

\section{Cluster Sensitivity for M08 and M09}
\ifLBResultsReady
For M08, the synchronized entity-cluster hierarchical interval is
$[\LBMEightClusterCILow,\LBMEightClusterCIHigh]$; each entity draw is shared by
all five seeds, models, and strengths in its dataset, with seed-specific effects
retained for outer resampling.  This is descriptive uncertainty, not an
equivalence or practical-null test.  For M09, the synchronized
designed-category reweighting interval is
$[\LBMNineReweightingLow,\LBMNineReweightingHigh]$; each source-category draw is
shared across models within dataset/seed/strength.  The latter is explicitly
descriptive because the eight categories exhaust a constructed mechanism and
are not a sample from a source population.  M09 is therefore only a
representation-conditional descriptive counterexample within this registry.
Both analyses use 200 inner draws per shared group and 5,000 outer draws over
tasks and seeds, with complete 2,500-cell coverage per mechanism.  Sensitivity
resampling changes uncertainty or reweighting, not the semantic C label.
\else
Cluster-sensitive estimates remain blocked pending complete prediction-file
coverage for all confirmatory M08/M09 cells.
\fi

\section{Detectability--Exploitability Diagnostics}
\ifLBResultsReady
The global mechanism-level Spearman correlation is \LBGlobalSpearman{} (95\%
joint-paired D--X bootstrap interval
$[\LBGlobalSpearmanCILow,\LBGlobalSpearmanCIHigh]$); identical dataset/seed draws
resample both axes.  Category-only and
category-plus-D in-sample $R^2$ values are \LBCategoryRSquared{} and
\LBCategoryPlusDRSquared{}, respectively; the within-category permutation
$p$-value for the increment is \LBIncrementalPermutationP{}.  LOMO $R^2$
values are \LBCategoryLomoRSquared{} and
\LBCategoryPlusDLomoRSquared{}.  These diagnostics are labeled
\emph{\LBRelationStatus{}} regardless of nominal significance because eleven
designed mechanisms do not establish out-of-registry transfer.
\else
Association diagnostics remain blocked and are not populated with pilot
values.
\fi

\section{Real-Data Boundary Audits}
The five fixed cases are UCI Bank Marketing \citep{uciBankMarketing}, a local
Lending Club accepted-loan mirror, U.S. flight delays \citep{btsOnTime},
Chicago food inspections \citep{chicagoFoodInspections}, and NYC 311
\citep{nyc311}.  The final lineage table
must bind the train-fitted-category preprocessing v2 protocol: category
vocabularies are learned from training rows, unseen validation/test categories
map to a reserved unknown value, and globally encoded date strings are dropped.
The earlier full-table-vocabulary outputs are superseded.  The lineage table
must include source filename and SHA-256, row-selection rule, chronological
sort field, prediction boundary, contaminated fields, sample size, split rule,
and all adapter exclusions.  Lending Club uses accepted applications only;
its post-origination servicing, last-credit-pull, hardship, and payment-plan
fields are post-boundary.  The flight-delay boundary is immediately before
scheduled departure: only an explicit schedule allowlist is legitimate, while
all actual departure, arrival, and diversion fields are post-event.  Bank
Marketing uses the canonical full dataset; no synthetic fallback is an
evaluable real-data row.  These declarations are adapter contracts rather than
labels inferred from the observed paired harm.

The Bank Marketing record is CC BY 4.0 and the canonical
\texttt{bank-additional-full} table contains 41,188 rows.  BTS, Chicago, and
NYC entries cite their official government portals.  The Lending mirror
\texttt{accepted\_2007\_to\_2018Q4} has no verified upstream provenance or
license in the current artifact: it is local-only and must not be redistributed
unless the authors supply and verify both.

\ifLBResultsReady
Across \LBNaturalTaskCount{} cases and \LBNaturalModelCount{} CPU learners,
mean paired harm is \LBNaturalMeanHarm{} with descriptive task-bootstrap
interval $[\LBNaturalCILow,\LBNaturalCIHigh]$ and exact task-level sign-flip
$p=\LBNaturalSignFlipP$.  All task effects positive: \LBNaturalAllPositive{}.
The evidence status is \emph{\LBNaturalStatus{}} and must not be rewritten as a
population-level, transfer, or detector-generalization result.
\else
Real-data numerical summaries remain blocked pending the same final claim
release as the controlled results.
\fi

\section{Negative, Superseded, and Out-of-Scope Evidence}
Legacy eight-model summaries, surrogate neural-model outputs, the pre-bundle
TabM partial run, old 10,083/4,032-cell ledgers, and earlier natural-transfer
tables are superseded evidence and cannot source this paper.  Corrected results
must live under the independent corrected-v2 namespace with checksums.  A
negative result for the pre-existing directional category claim remains
reportable as \texttt{NOT\_SUPPORTED}; a
\texttt{PENDING}, \texttt{REFUTED}, \texttt{INTEGRITY\_HOLD}, or incomplete
claim cannot generate a paper macro.

\paragraph{Frozen M10 protocol amendment.}
The original M10 strict-view constructor incorrectly removed the injected
legitimate duplicate together with the fields marked by
\texttt{leakage\_mask}; the strict and permissive views therefore did not share
the same legitimate base.  The frozen amendment defines the strict matrix as
\texttt{task.X[:, \string~leakage\_mask]}.  The original 2,000 CPU and 500 TabM
M10 rows are superseded and replaced under this corrected pairing.  Canonical
validation excludes every original M10 run identifier and still requires
exactly 27,500 model-training cells.  Superseded M10 rows cannot source a claim,
table, or figure.

Operational metadata and governance modules are engineering/audit work outside
the main confirmatory claim set.  Their omission is not evidence for or against
their efficacy.

\section{Claim--Evidence Matrix}
The final release should map each main-paper statement to (i) a stable claim
identifier, (ii) its scope and, only for the directional category contrast, its
support rule, (iii) final status, (iv) exact
metric fields, (v) canonical source hashes, and (vi) allowed and prohibited
wording.  The paper generator consumes only the six controlled claim IDs and
the case-study-only natural block.  Every other state is excluded by design.

\section{Complete Machine-Generated Tables}
\ifLBResultsReady
\IfFileExists{generated/result_tables.tex}{%
\input{generated/result_tables.tex}}{%
\PackageError{LeakBench-Tab}{Missing final result tables}{Run
source_data/generate_result_tables.py from the final claim release.}}
\else
Complete registry, mechanism, model, diagnostic, strength, natural-case, and
claim-scope tables are blocked with the numerical results.  Pilot tables are
not rendered in their place.
\fi

\section{Reproducibility Checklist and Ledger}
The archival package should provide the frozen corrected protocol, task-bundle
manifest, canonical 27,500-cell table, raw failure-preserving CPU and official
TabM tables, retained M08/M09 predictions, the raw and amended canonical
diagnostic tables, the diagnostic RNG amendment freeze and impact audit,
the statistical-amendment freeze and input/output hash manifests,
natural-task lineage, all statistics tables, \texttt{paper\_claims.json},
generated macro source hash, environment
locks, and a release validator.  Reproduction commands must distinguish CPU
and official GPU runs and must never regenerate task bundles in a different
dependency environment.

\section{Limitations and Ethical Considerations}
The benchmark uses synthetic controlled tasks, designed mechanisms, binary
classification, and a limited model set.  The real-data component contains five
selected boundary audits rather than a random task sample.  Leakage findings
can reveal risky data practices; releases should avoid exposing credentials,
personal information, or restricted raw data.  Task manifests and derived
statistics should be public only when their source licenses permit it.

\bibliography{references}
\end{document}
